\section{Hash Tables as Non-Sequential Component Access Structures}

Sequential containers teach one model of storage: components arranged in a visible order and accessed by integer position. A string, list, or tuple can answer questions such as ``what is at index \texttt{0}?'' or ``what slice lies between these two boundaries?'' That model is useful, but it is not the only way to organize components.

Sets and dictionaries use a different strategy. They are not primarily built around position. They are built around hash-based access. A set asks whether a value is present. A dictionary asks which value is associated with a key. In both cases, Python uses a hash value to move quickly toward the internal table area where the element or key may be stored.

This section introduces that change of perspective. The focus is not yet on all set and dictionary operations, but on the access architecture underneath them: how Python moves from sequential indexing to hash-derived placement, why equality is still needed after hashing, and why only hash-stable objects may safely participate as set elements or dictionary keys.